\chapter*{Notation}


I use the following conventional notation in this book.


\subsection*{Math}

\begin{tabular}{ll}
\hline 
$\binom{n}{m}$ & ``$n$ choose $m$'' which equals $\frac{n!}{m!(n-m)!}$ \\
$\sum$ & summation, e.g., $\sumn a_i = a_1 + \cdots + a_n$ \\
$I(\cdot)$ & indicator function, i.e.,  $I(A)=1$ if $A$ happens and 0 otherwise\\
$\#$ & counting the number of units in a set\\
$\approx$ & approximately equal \\
$\propto$ & proportional to (by dropping some unimportant constant) \\
$\text{logit}$ &  $\text{logit}(x) = \log \frac{x}{1-x}$ \\
$\expit$ & $\expit(x) =   \frac{e^x}{1+e^x} =  (1 + e^{-x})^{-1}$ \\
\hline 
\end{tabular}

\subsection*{Basic probability and statistics}

\begin{tabular}{ll}
\hline 
$\pr(\cdot)$ & probability \\
$E(\cdot)$ & expectation of a random variable\\
$\var(\cdot)$ & variance of a random variable \\
$\cov(\cdot)$ & covariance between random variables\\
$\asim$ & ``$A \asim B$'' means that $A$ and $B$ have the same asymptotic distribution\\
$\rho_{YX}$ & Pearson correlation coefficient between $Y$ and $X$ \\
$R^2_{YX}$ & squared multiple correlation coefficient between $Y$ and $X$ \\
\hline 
\end{tabular}


\subsection*{Random variables}

\begin{tabular}{ll}
\hline 
Bernoulli$(p)$ & Bernoulli distribution with probability $p$\\
Binomial$(n,p)$& Binomial distribution with $n$ trials and probability $p$\\
$\textsc{N}(\mu, \sigma^2)$ & Normal distribution with mean $\mu$ and variance $\sigma^2$\\
$t_\nu$ & $t$ distribution with degrees of freedom $\nu$ \\
$\chi^2_\nu$ & chi-squared distribution with degrees of freedom $\nu$\\
$z_{1-\alpha/2}$ & the $1-\alpha/2$ upper quantile of $\N01$, e.g., $z_{0.975}=1.96$\\
\hline 
\end{tabular}


\subsection*{Causal inference}

\begin{tabular}{ll}
\hline 
 $Y_i(1), Y_i(0)$ & potential outcomes of unit $i$ under treatment and control\\
$\tau$ & finite-population average causal effect $\tau = n^{-1}\sumn \{Y_i(1)-Y_i(0)\}$ \\
$\tau$ & super-population average causal effect $\tau = E \{Y(1)-Y(0)\}$\\
$\mu_1(X)$ & outcome model $\mu_1(X) = E(Y\mid Z=1, X)$\\
$\mu_0(X)$ & outcome model $\mu_0(X) = E(Y\mid Z=0, X)$\\
$e(X)$ & propensity score $e(X) = \pr(Z=1\mid X)$ \\
$\tau_\cp$ & complier average causal effect $ \tau_\cp = E\{ Y(1) - Y(0)\mid U=\cp \} $ \\
$U$ & unmeasured confounder \\
$U$ & latent compliance status $U=(D(1), D(0))$ \\
$Y(z,M_{z'})$ & nested potential outcome (for mediation analysis)\\
 \hline 
\end{tabular}